\section{The exact symbol calculus and the established kernel theorem}
\label{sec:symbols}

\subsection{Rationalized exterior symbols}

For a number field $K$, put
\[
 \A(K)=K^\times\otimes_\Z\Q,
 \qquad \cl{uv}=\cl u+\cl v.
\]
All exterior powers in this article are over $\Q$. Roots of unity vanish
in $\A(K)$. Embeddings of fields induce injective maps on these vector
spaces and on their exterior powers, so we identify old-field symbols
with their images in larger fields.

\begin{definition}[Formal reductions]\label{def:formal}
Let $\Pcal(K)$ be the rationalized pre-Bloch group: the vector space on
symbols $[z]$, modulo five-term relations, with $[0]=[1]=0$. Write
\[
 \partial[z]=\cl z\wedge\cl{1-z}\quad(z\ne0,1),
 \qquad \B(K)=\ker\partial.
\]
For a Galois-invariant element $\eta\in\Pcal(K)$, a
\emph{formal logarithmic reduction} means $\eta=0$. A
\emph{formal rational-dilogarithmic reduction} means that $\eta$ belongs
to the image of $\Pcal(\Q)$.
\end{definition}

The terminology suppresses products of logarithms of algebraic numbers
and rational multiples of $\pi^2$ in evaluation. Compatible complex
logarithms are allowed in this definition. The worked analytic identities
below use only real logarithms of positive algebraic numbers and
$\Li_2$ on $(0,1)$.

For finite Galois number fields, the standard rationalized inputs are
\begin{equation}\label{eq:K-inputs}
 \B(K)^{\Gal(K/\Q)}=\B(\Q)=0,
 \qquad \partial\Pcal(\Q)=\Lambda^2\A(\Q).
\end{equation}
These are Bloch-group descent, the torsion of $\B(\Q)$ and the torsion
of $K_2(\Q)$; the last assertion uses the identification of the
Steinberg quotient with $K_2$ for a field. They are the same inputs used
in \cite[\S7.2]{RZ}. For the Bloch-group and $K$-theoretic framework,
see \cite{ZagierDilog,CGZ}. Consequently, for invariant $\eta$,
\begin{align}
 \eta=0&\ \Longleftrightarrow\ \partial\eta=0,
 \label{eq:zero-detection}\\
 \eta\in\operatorname{im}\Pcal(\Q)
 &\ \Longleftrightarrow\
 \partial\eta\in\Lambda^2\A(\Q).
 \label{eq:rational-detection}
\end{align}
For the reverse implication in the second line, choose a rational
pre-Bloch element with the same boundary and apply the first line to
the invariant difference. This makes clear where the established
$K$-theory input enters; it is not replaced by a numerical observation.

\subsection{Cyclotomic symbols and the direct boundary formula}

For $m\ge3$ and $(a,m)=1$, define
\begin{equation}\label{eq:beta}
 \beta_m(a)=\sum_{k=1}^{m-1}
 \cl{1-\zeta_m^{ak}}\wedge\cl{1-\zeta_m^k},
 \qquad \zeta_m=e^{2\pi i/m}.
\end{equation}
Use $\beta_1=\beta_2=0$. The excluded term $k=0$ is important: there is
no multiplicative-group element $\cl0$. Reindexing gives
\begin{equation}\label{eq:symmetries}
 \beta_m(-a)=\beta_m(a),\qquad
 \beta_m(a\inv)=-\beta_m(a).
\end{equation}
Every $\beta_m(a)$ is Galois invariant. The inverse is taken modulo $m$.

For coprime positive $p,q$, the finite formula of \cite[Theorem 3]{RZ}
defines an invariant element $\xi_{p/q}\in\Pcal(\Q(\zeta_{pq}))$
representing the dilogarithmic part of $F(p/q)-F(1)$. Explicitly it is
\begin{equation}\label{eq:xi-finite}
 \xi_{p/q}=\sum_{\alpha^p=1}
 \left(\sum_{\substack{\beta^q=1\\\beta\ne1}}h(\alpha,\beta)
       -\sum_{\substack{\beta^p=1\\\beta\ne1}}h(\alpha,\beta)\right),
 \quad
 h(\alpha,\beta)=
 \left[\frac{\beta}{\beta-1}\right]
 -\left[\frac{\alpha-\beta}{1-\beta}\right].
\end{equation}

\begin{proposition}[Established boundary formula]\label{prop:boundary}
For coprime $p,q>0$,
\begin{equation}\label{eq:boundary-F}
 \partial\xi_{p/q}
 =\beta_q(p)-\beta_p(q)-\cl p\wedge\cl q.
\end{equation}
\end{proposition}
\begin{proof}
This is the boundary used in both \cite{RZ} and \cite{Repo}; we recall its
direct verification to avoid relying on the assertion audited in
Section~\ref{sec:audit}. Terms with $\alpha=1$ or $\alpha=\beta$
have zero boundary. For the others set
$A=\cl{1-\alpha}$, $B=\cl{1-\beta}$, $C=\cl{\alpha-\beta}$.
Since roots of unity vanish, direct expansion gives
\[
 \partial h(\alpha,\beta)=-C\wedge A+C\wedge B+B\wedge A.
\]
The subtracted double sum in \eqref{eq:xi-finite} has zero boundary by
antisymmetry under $\alpha\leftrightarrow\beta$. In the first double
sum, the product identities
\[
 \prod_{\substack{\beta^q=1\\\beta\ne1}}(\alpha-\beta)
 =\frac{\alpha^q-1}{\alpha-1},\qquad
 \prod_{\substack{\alpha^p=1\\\alpha\ne1}}(\beta-\alpha)
 =\frac{\beta^p-1}{\beta-1}
\]
make the first two wedge contributions $-\beta_p(q)$ and $\beta_q(p)$.
The last contribution is $\cl q\wedge\cl p$, since the products of
$1-\alpha$ and $1-\beta$ over the nontrivial roots are $p$ and $q$.
Coprimality ensures that no additional common root was omitted.
\end{proof}

\subsection{A faithful model on the symbol span}

Put $G_m=(\Z/m\Z)^\times/\{\pm1\}$ for $m\ge3$. The following theorem
is an existing result of the manuscript \cite{Repo}; its proof is
included because the new layer argument uses its Gram matrix, not just
its conclusion.

\begin{theorem}[Existing all-conductor kernel theorem]\label{thm:kernel}
For $c_a\in\Q$,
\begin{equation}\label{eq:kernel}
 \sum_{a\in G_m}c_a\beta_m(a)=0
 \quad\Longleftrightarrow\quad c_a=c_{a\inv}\quad(a\in G_m).
\end{equation}
Thus the map on the symbol span
\begin{equation}\label{eq:group-model}
 \beta_m(a)\longmapsto \ee_a-\ee_{a\inv}\in\Q[G_m]
\end{equation}
is faithful, and
\begin{equation}\label{eq:beta-zero}
 \beta_m(a)=0\quad\Longleftrightarrow\quad
 a^2\equiv\pm1\pmod m.
\end{equation}
\end{theorem}
\begin{proof}
For $1\le k<m$ define vectors and a matrix
\[
 v_k=(\log|1-\zeta_m^{bk}|)_{b\in G_m},
 \qquad A_m=\sum_{k=1}^{m-1}v_kv_k^{\mathsf T}.
\]
Let $(P_av)_b=v_{ba}$. Then $P_av_k=v_{ak}$, so $A_m$ commutes with
all $P_a$.

The vectors $v_k$ span the regular representation. Indeed their span
is invariant, and it has a nonzero projection to each character. For a
nontrivial even character of primitive conductor $f\mid m$, use
$k=m/f$. Its Fourier coefficient is a positive integer multiple of
\[
 \sum_{b\in G_f}\chi^*(b)\log|1-\zeta_f^b|
 =-\tfrac12\tau(\chi^*)L(1,\overline{\chi^*})\ne0.
\]
The equality follows by Abel-regularizing
$\log|1-r e^{i\theta}|=-\sum_{n\ge1}r^n\cos(n\theta)/n$ and applying
the primitive Gauss-sum identity. Nonvanishing uses the nonzero primitive
Gauss sum and Dirichlet's theorem $L(1,\chi)\ne0$ for nonprincipal
characters \cite[25.15.9]{DLMF}. The trivial character occurs using
$k=m/\ell$ for a prime divisor $\ell$ of $m$: the logarithmic product
is $\log\ell$; for $\ell=2$ the vector is the nonzero constant vector
$\log2$. All nontrivial roots, not just primitive $m$th roots, are used.
Therefore $A_m$ is positive definite.

Apply the exterior square of the logarithm map
$\A(\Q(\zeta_m))\to\R^{G_m}$, and identify $v\wedge w$ with
$vw^{\mathsf T}-wv^{\mathsf T}$. Its value on a symbol is
\begin{equation}\label{eq:skew-image}
 \ell_\wedge(\beta_m(a))
 =P_aA_m-A_mP_a^{\mathsf T}
 =A_m(P_a-P_{a\inv}).
\end{equation}
A zero combination of symbols therefore gives
$\sum_a(c_a-c_{a\inv})P_a=0$ after multiplication by $A_m^{-1}$.
The regular permutation matrices are linearly independent, so the
coefficients are symmetric. Conversely symmetric coefficients cancel
by \eqref{eq:symmetries}; self-inverse classes contribute zero over
$\Q$. This proves all assertions.
\end{proof}

\begin{remark}
The exterior logarithm map is not asserted to be injective on all of
$\Lambda^2\A(K)$. Its explicitly computed value suffices on the
particular span in Theorem~\ref{thm:kernel}. Positivity above is proved
by character nonvanishing, not by a floating-point determinant.
\end{remark}

\subsection{Separating coprime conductors}

\begin{lemma}[Normalized norm separation]\label{lem:norm}
Let $(r,s)=1$, with $U$ a rational combination of the $\beta_r(a)$ and
$V$ a rational combination of the $\beta_s(b)$, viewed in a common
cyclotomic field. Then
\[
 U-V\in\Lambda^2\A(\Q)\quad\Longleftrightarrow\quad U=V=0.
\]
Moreover, normalized norm to $\Q$ kills each such symbol combination.
\end{lemma}
\begin{proof}
The linear normalized field norm to $E$ is
$N_E\cl u=\cl{\Norm_{K/E}u}/[K:E]$. It fixes old-field elements;
apply its exterior square to wedges. Coprime cyclotomic fields are
linearly disjoint over $\Q$. The norm of either factor of a term of
$\beta_s(b)$ to $\Q(\zeta_r)$ is the same rational vector, since the
factors are Galois conjugates in $\Q(\zeta_s)$. Thus that wedge is
killed, while every old $r$-conductor wedge is fixed. Reversing the
roles gives the other projection. If $U-V=W$ is rational, the two
projections give $U=W$ and $-V=W$; projecting either again to $\Q$
gives $W=0$. The cases of conductor $1$ or $2$ follow from the
zero-symbol conventions. The same conjugate-factor argument proves
the last assertion.
\end{proof}
